\documentclass[10pt,a4paper]{amsart}
\usepackage[margin=19mm]{geometry}
\usepackage{amsmath,amssymb,mathtools}
\usepackage[scr=rsfs,cal=euler]{mathalfa}
\usepackage{xfrac}
\usepackage[unicode,hidelinks,bookmarks=false]{hyperref}
\newcommand{\ie}{i.e.\ }
\newcommand{\cf}{cf.\ }
\newcommand{\resp}{respectively\ }
\newcommand{\Subsection}{\subsection}
\providecommand{\nobreakdash}{\nobreak\hbox{-}}
\setlength{\emergencystretch}{2em}
\multlinegap0pt
\usepackage{bm}
\usepackage{bbm}
\usepackage{dsfont}
\usepackage{xspace}
\usepackage{cancel}
\usepackage{mathtools}
\usepackage[shortlabels]{enumitem}
\usepackage{amsthm}
\makeatletter
\@ifundefined{lem}{\newtheorem{lem}{Lemma}}{}
\makeatother
\newcommand{\R}{{\mathbb{R}}}
\title[Optimal regularity results in Sobolev-Lorentz spaces for lin]{Optimal regularity results in Sobolev-Lorentz spaces for linear elliptic equations with $L^1$- or measure data}
\author{Hyunseok Kim, Young-Ran Lee, Jihoon Ok}
\date{Source: arXiv:2506.15005v2 (CC BY 4.0)}
\begin{document}
\maketitle

\noindent\emph{Source and licence.} Optimal regularity results in Sobolev-Lorentz spaces for linear elliptic equations with $L^1$- or measure data. Version arXiv:2506.15005v2 (CC BY 4.0), distributed under the Creative Commons Attribution 4.0 International licence, as recorded in the arXiv licence metadata for this exact version and shown on the abstract page https://arxiv.org/abs/2506.15005v2. Full rights notice: licenses/arxiv-2506.15005-CC-BY-4.0.txt.\par\smallskip

\noindent\textit{Editorial scope.} Counted unit: the Lem whose source label is \texttt{properties of SL on domains} in the article (source0.tex lines 771--778, source label \texttt{properties of SL on domains}), delivered together with the article's own complete original proof of it (source0.tex lines 781--864, one \texttt{proof} environment of 84 lines). No mathematics is written, repaired or re-derived: statement and proof are the article's own text line for line. Required context retained uncounted: widehat-L (lines 472--475); properties of SL on Rn (lines 521--529); complex interp for S (lines 581--583); interp for SL on domains (lines 599--609); comp-interp-S on domains (lines 665--667); real-interp-S on domains (lines 672--685). Every definition, display and label of those lines is retained unchanged. Omitted from this package and not redistributed: 1--471, 476--520, 530--580, 584--598, 610--664, 668--671, 686--770, 779--780, 865--1946. Nothing inside a retained excerpt is omitted or altered: every retained body is delivered exactly as the article's lines are written, so no omission receipt of any kind accompanies this package. Citations: the retained excerpts carry no citation command at all, so no bibliography entry of the article is transcribed and no reference is redistributed. Reference adaptation: none was needed and none was made; every reference inside the delivered unit resolves inside it, so no unresolved reference or citation is produced. Printed slips kept literal: no printed word, symbol, hypothesis or number of the counted statement or of its proof was altered. Layout and class: the article's own class is \texttt{amsart} with packages including amsthm, amsmath, amsfonts, geometry, amssymb, theoremref, mathrsfs, enumerate, hyperref, xcolor, kotex; the wrapper is the lane's pinned \texttt{amsart} editorial wrapper at 10pt on A4, which restates the theorem-like environments the retained text needs and adds the article's own macro definitions that the retained text needs (\texttt{\textbackslash R}), with extra wrapper packages (bm, bbm, dsfont, xspace, cancel, mathtools, enumitem). The delivered numbering is the unit's own. Assets: no retained span contains a figure environment, a graphics-inclusion command, a PDF-page inclusion, a table or a TikZ picture; no image file is copied into this package, the package declares no asset dependency and no image file is redistributed with it. Licence: version arXiv:2506.15005v2 of this article is distributed under the Creative Commons Attribution 4.0 International licence, as recorded in the arXiv licence metadata for this exact version and shown on the abstract page https://arxiv.org/abs/2506.15005v2. A full rights notice is delivered at licenses/arxiv-2506.15005-CC-BY-4.0.txt. Characters: every retained excerpt is pure ASCII, so the delivered engine, XeLaTeX with its default Latin Modern text font, reports no missing character.
\begin{equation}\label{widehat-L}
\left[ \widetilde{L}^{p,q} (\Omega )\right]^*
= L^{p',q'} (\Omega )  \quad\mbox{for}\,\, 1<p<\infty, \,\,1 \le q \le \infty.
\end{equation}

\begin{lem}\label{properties of SL on Rn} Let $ -\infty<  \alpha, \alpha_0 < \infty$,  $1< p , p_0< \infty$, and    $1 \le q, q_0 \le \infty$.
\begin{enumerate}[{\upshape (i)}]
\item  If $\alpha  \ge \alpha_0$  and $q \le q_0 $, then   $L_\alpha^{p,q}(\R^n ) \hookrightarrow L_{\alpha_0}^{p,q_0 }(\R^n )  $. In particular, if $\alpha > 0$, then $L_\alpha^{p,q}(\R^n ) \hookrightarrow L^{p,q  }(\R^n )\hookrightarrow L_{-\alpha}^{p,q}(\R^n )$.
\item  If $\alpha-\alpha_0 = n/p-n/p_0 >0$, then $L_\alpha^{p,q}(\R^n)\hookrightarrow L_{\alpha_0}^{p_0 , q}(\R^n)$. In particular, if $0<  \alpha p < n $, then $L_\alpha^{p,q}(\R^n)\hookrightarrow L^{np/(n-\alpha p),q}(\R^n)$.
\item   If $\alpha p=n$, then $L_\alpha^{p,1}(\R^n)\hookrightarrow C_0 (\R^n )$, where $C_0 (\R^n )$ is the closure of $C_c^\infty (\R^n)$ in $L^\infty (\R^n )$.
\item   If $k$ is a positive integer, then $L_k^{p,q}(\R^n) =W^{k,p,q}(\R^n )$, where $W^{k,p,q}(\R^n )$ is the Banach space  of  all $f \in L^{p,q}(\R^n )$ such that $D^\gamma f \in L^{p,q}(\R^n )$ for all multi-indices $\gamma$ with $|\gamma | \le k$.
\item  If $q < \infty$, then $C_c^\infty (\R^n)$ is dense in  $L_\alpha^{p,q}(\R^n)$ and $\left[L_\alpha^{p,q}(\R^n) \right]^* = L_{-\alpha}^{p',q'}(\R^n)$.
\end{enumerate}
\end{lem}

\begin{equation}\label{complex interp for S}
L_{\alpha}^{p} (\Omega )  = \left[ L_{\alpha_0}^{p_0} (\Omega ) ,L_{\alpha_1}^{p_1} (\Omega ) \right]_{\theta}  .
\end{equation}

\begin{lem}\label{interp for SL on domains} Let $0 \le   \alpha  <\infty$, $1< p <\infty$, and $ 1\le q  \le \infty$.
\begin{enumerate}[{\upshape (i)}]
\item $E_\Omega$ is bounded from $L_\alpha^{p,q}(\Omega )$ into $L_\alpha^{p,q}(\R^n )$.
\item  If      $1< p_0 \neq p_1 < \infty$, $0<\theta < 1$, and $1/p=(1-\theta)/p_0 +\theta/p_1$,    then
    $$
    L_\alpha^{p,q}(\Omega ) = \left(\,  L_\alpha^{p_0}(\Omega ) , L_\alpha^{p_1}(\Omega ) \, \right)_{\theta,q}   .
    $$
\item     If $k$ is a positive integer, then $L_k^{p,q}(\Omega ) =W^{k,p,q}(\Omega )$.
\item  If   $q<\infty$, then  $C^\infty (\overline{\Omega}) $ is dense in $L_\alpha^{p,q}(\Omega )$, where $C^\infty (\overline{\Omega}) =R_\Omega \left(C_c^\infty (\R^n ) \right)$.
\end{enumerate}
\end{lem}

\begin{equation}\label{comp-interp-S on domains}
\widetilde{L}_{\alpha}^{p} (\Omega )  = \left[\widetilde{L}_{\alpha_0}^{p_0} (\Omega ) ,\widetilde{L}_{\alpha_1}^{p_1} (\Omega ) \right]_{\theta} \quad\mbox{and}\quad L_{-\alpha}^{p} (\Omega )  =   \left[ L_{-\alpha_0}^{p_0} (\Omega ) ,L_{-\alpha_1}^{p_1} (\Omega ) \right]_{\theta}.
\end{equation}

\begin{lem}\label{real-interp-S on domains} Let $0 \le   \alpha  <\infty$, $1< p  < \infty$,  and $ 1\le q  \le \infty$.
\begin{enumerate}[{\upshape (i)}]
\item  If      $1< p_0 \neq p_1 < \infty$, $0<\theta < 1$, and $1/p=(1-\theta)/p_0 +\theta/p_1$,    then
    \[
 R_\Omega  (\widetilde{L}_\alpha^{p,q} (\overline{\Omega} ) )  = \left( \widetilde{L}_\alpha^{p_0}(\Omega ) ,\widetilde{L}_\alpha^{p_1}(\Omega ) \right)_{\theta,q}
\quad\mbox{and}\quad
     L_{-\alpha}^{p,q} (\Omega )   = \left( L_{-\alpha}^{p_0} (\Omega ) ,L_{-\alpha}^{p_1} (\Omega ) \right)_{\theta,q}.
\]
\item If   $q<\infty$, then  $C_c^\infty (\Omega) $ is dense in $R_\Omega  (\widetilde{L}_\alpha^{p,q} (\overline{\Omega} ) )$, that is,   $\widetilde{L}_\alpha^{p,q}(\Omega )= R_\Omega  (\widetilde{L}_\alpha^{p,q} (\overline{\Omega} ) )$.
\item If  $k$ is  a positive integer, then $\widetilde{L}_k^{p,q}(\Omega ) $ is the closure of $C_c^\infty (\Omega) $ in $L_k^{p,q}(\Omega )$.
\item If $q>1$, then $C_c^\infty (\Omega )$ is dense in $L_{-\alpha}^{p',q'} (\Omega ) $ and $ \left[ L_{-\alpha}^{p', q'}(\Omega ) \right]^* = R_\Omega  (\widetilde{L}_\alpha^{p,q} (\overline{\Omega} ) )$.
    %$ L_{\alpha}^{p,q}(\Omega )$.
\end{enumerate}
\end{lem}

\begin{lem}\label{properties of SL on domains} Let $-\infty < \alpha, \alpha_0 <\infty$, $1< p, p_0 <\infty$, and $ 1\le q ,q_0 \le \infty$.
\begin{enumerate}[{\upshape (i)}]
\item  If $\alpha  \ge \alpha_0$  and $q \le q_0 $, then   $L_\alpha^{p,q}(\Omega ) \hookrightarrow L_{\alpha_0}^{p,q_0 }(\Omega ) $. In particular, if $\alpha > 0$, then $L_\alpha^{p,q}(\Omega ) \hookrightarrow L^{p,q  }(\Omega )\hookrightarrow L_{-\alpha}^{p,q}(\Omega )$.
\item  If  $p >  p_0$,  then    $L_\alpha^{p,q}(\Omega ) \hookrightarrow L_{\alpha}^{p_0,q_0 }(\Omega ) $.
\item    If $\alpha-\alpha_0 = n/p-n/p_0 >0$, then $L_\alpha^{p,q}(\Omega )\hookrightarrow L_{\alpha_0}^{p_0 , q}(\Omega )$. In particular,  if $0<  \alpha p < n $, then $L_\alpha^{p,q}(\Omega)\hookrightarrow L^{np/(n-\alpha p),q}(\Omega)$.
\item     If $\alpha p=n$, then $L_\alpha^{p,1}(\Omega)\hookrightarrow C  (\overline{\Omega} )$ and $\widetilde{L}_\alpha^{p,1}(\Omega)\hookrightarrow C_0  (\Omega )$, where $C_0 (\Omega )$ is the closure of $C_c^\infty (\Omega )$ in $L^\infty (\Omega)$
\end{enumerate}
\end{lem}

\begin{proof} To prove Part (i), suppose that $\alpha  \ge \alpha_0$  and $q \le q_0 $. If $\alpha_0 \ge 0$, then the embeddings   $ L_\alpha^{p,q}(\Omega )  \hookrightarrow    L_{\alpha_0}^{p,q_0 }(\Omega )\hookrightarrow    L^{p,q_0 }(\Omega )$  immediately follow    from Lemma \ref{interp for SL on domains} (i) and Lemma \ref{properties of SL on Rn} (i).
Suppose that $\alpha_0 < 0$. Then since $-\alpha_0 >  0$ and $q_0 '\le q'$, it follows that  $ L_{-\alpha_0}^{p',q_0'}(\R^n )  \hookrightarrow    L^{p',q '}(\R^n ) $. Using this embedding, we easily deduce that
$$
\widetilde{L}_{-\alpha_0}^{p',q_0 '}(\Omega )     \hookrightarrow  \widetilde{L}^{p',q'}(\Omega ).
$$
%On the other hand, since $L^{p',q'} (\Omega )= \left(\, L^{p_1 '}(\Omega ), L^{p_2 '}(\Omega ) \, \right)_{1/2 , q'}$  for some $p_1 , p_2$ with   $2/p= 1/p_1 +1/p_2$, it follows from  the duality theorem (see, e.g., \cite[Theorem 3.7.1 and its Remark]{bergh}) that
%\[
%L^{p,q}(\Omega )  = \left(\, L^{p_1  }(\Omega ), L^{p_2  }(\Omega ) \, \right)_{1/2 , q} = \left(\, L^{p_1 '  }(\Omega )^* , L^{p_2 ' }(\Omega )^*  \, \right)_{1/2 , q} = \left[  \widetilde{L}^{p',q'}(\Omega ) \right]^* .
%\]
%Therefore,
%\[
%L^{p,q}(\Omega )  = \left[ \widetilde{L}^{p',q'}(\Omega ) \right]^* \hookrightarrow \left[ \widetilde{L}_{-\alpha_0}^{p',q_0 '}(\Omega ) \right]^* = L_{\alpha_0}^{p,q_0 }(\Omega )   .
%\]
Hence by the duality result  (\ref{widehat-L}),
\[
L^{p,q}(\Omega )  = \left[ \widetilde{L}^{p',q'}(\Omega ) \right]^* \hookrightarrow \left[ \widetilde{L}_{-\alpha_0}^{p',q_0 '}(\Omega ) \right]^* = L_{\alpha_0}^{p,q_0 }(\Omega )   .
\]
Therefore, if $\alpha \ge 0 > \alpha_0$, then
$L_{\alpha}^{p,q  }(\Omega ) \hookrightarrow  L^{p,q}(\Omega )   \hookrightarrow L_{\alpha_0}^{p,q_0 }(\Omega )$.
Suppose finally that $\alpha < 0$.  Then since $-\alpha_0 \ge -\alpha > 0$ and $q_0 '\le q'$, we have
$$
L_{-\alpha_0}^{p',q_0'}(\R^n )  \hookrightarrow    L_{-\alpha}^{p',q '}(\R^n )
\quad\mbox{and so}\quad
\widetilde{L}_{-\alpha_0}^{p',q_0 '}(\Omega )  \hookrightarrow  \widetilde{L}_{-\alpha}^{p',q  '}(\Omega )   .
$$
Therefore,
\[
L_\alpha^{p,q}(\Omega )  = \left[ \widetilde{L}_{-\alpha}^{p',q'}(\Omega ) \right]^* \hookrightarrow \left[ \widetilde{L}_{-\alpha_0}^{p',q_0 '}(\Omega ) \right]^* = L_{\alpha_0}^{p,q_0 }(\Omega )  .
\]
This proves Part (i).

To prove Part (ii), suppose that $p>p_0$.  The case $\alpha =0$ is trivial. Suppose that $\alpha >0$. If $k$ is a positive  integer, then $L_k^{p,p}(\Omega ) =W^{k,p}(\Omega )  \hookrightarrow  W^{k,p_0}(\Omega )= L_k^{p_0 ,p_0}(\Omega )$.
Moreover, it follows from Lemma \ref{real-interp-S on domains} (iii) that $\widetilde{L}_k^{p,p}(\Omega ) \hookrightarrow \widetilde{L}_k^{p_0 ,p_0}(\Omega )$.
Hence, if $k $ is any integer such that   $k >  \alpha $, then by   (\ref{complex interp for S}) and (\ref{comp-interp-S on domains}), we   have
$$
 L_\alpha^{p,p} (\Omega )  = \left[ L^p (\Omega ), L_k^{p}(\Omega )\right]_{\alpha /k} \hookrightarrow \left[ L^{p_0} (\Omega ), L_k^{p_0}(\Omega )\right]_{\alpha /k}=L_\alpha^{p_0 , p_0} (\Omega )
$$
and similarly
$$
\widetilde{L}_\alpha^{p,p}(\Omega ) \hookrightarrow \widetilde{L}_\alpha^{p_0 ,p_0}(\Omega ).
$$
For the general case when $q \neq p$ or $q_0 \neq p_0$, we choose $1<p_1 , p_2 , p_3 <\infty$ such that  $p_1 < p_0 < p_2 < p <p_3$, $2/p_0 =1/p_1 +1/p_2$, and $2/p =1/p_2 +1/p_3$.
Then since $L_\alpha^{p_3}(\Omega )  \hookrightarrow  L_\alpha^{p_2}(\Omega )\hookrightarrow  L_\alpha^{p_1}(\Omega )$, it follows from  Lemma \ref{interp for SL on domains} (ii) that
\[
L_\alpha^{p,q}(\Omega ) = \left(\,  L_\alpha^{p_2}(\Omega ) , L_\alpha^{p_3}(\Omega ) \, \right)_{1/2,q}  \hookrightarrow  L_\alpha^{p_2}(\Omega ) \hookrightarrow   \left(\,  L_\alpha^{p_1}(\Omega ) , L_\alpha^{p_2}(\Omega ) \, \right)_{1/2,q_0} = L_\alpha^{p_0 ,q_0 }(\Omega ).
\]
Similarly, since $\widetilde{L}_\alpha^{p_3}(\Omega )  \hookrightarrow  \widetilde{L}_\alpha^{p_2}(\Omega )\hookrightarrow  \widetilde{L}_\alpha^{p_1}(\Omega )$, it follows from  Lemma \ref{real-interp-S on domains} (i) that
\[
R_\Omega  (\widetilde{L}_\alpha^{p,q} (\overline{\Omega} ) ) \hookrightarrow  R_\Omega  (\widetilde{L}_\alpha^{p_0 ,q_0} (\overline{\Omega} ) )
\]
and hence
\[
\widetilde{L}_\alpha^{p,q}(\Omega )   \hookrightarrow \widetilde{L}_\alpha^{p_0 ,q_0 }(\Omega ).
\]

Suppose next that  $\alpha < 0$. Then since $\widetilde{L}_{-\alpha}^{p_0 ',q_0 '}(\Omega )   \hookrightarrow \widetilde{L}_{-\alpha}^{p' ,q ' }(\Omega )$ and $\widetilde{L}_{-\alpha}^{p_0 ',q_0 '}(\Omega ) $ is dense in $ \widetilde{L}_{-\alpha}^{p' ,q ' }(\Omega )$, we immediately obtain
\[
L_\alpha^{p,q}(\Omega )  = \left[ \widetilde{L}_{-\alpha}^{p',q'}(\Omega ) \right]^* \hookrightarrow \left[ \widetilde{L}_{-\alpha}^{p_0 ',q_0 '}(\Omega ) \right]^* = L_{\alpha}^{p_0 ,q_0 }(\Omega )  ,
\]
which proves Part (ii).

To prove Part (iii), suppose  that $\alpha-\alpha_0 = n/p-n/p_0 >0$. If $\alpha_0 \ge 0$, then the embeddings   $ L_\alpha^{p,q}(\Omega )  \hookrightarrow    L_{\alpha_0}^{p_0,q }(\Omega )
%\hookrightarrow    L^{p_0,q }(\Omega )
$  immediately follow   from Lemma \ref{interp for SL on domains} (i) and Lemma \ref{properties of SL on Rn} (ii). Suppose   that $\alpha \le 0$. Then since $(-\alpha_0) -( -\alpha) = n/p_0' - n/p' > 0$, we have
 $$
L_{-\alpha_0}^{p_0',q'}(\R^n )  \hookrightarrow    L_{-\alpha}^{p',q'}(\R^n ) ,
%\hookrightarrow    L^{p',q '}(\Omega ),
$$
from which we easily deduce  that
\[
\widetilde{L}_{-\alpha_0}^{p_0',q'}(\Omega )  \hookrightarrow  \widetilde{L}_{-\alpha}^{p',q  '}(\Omega )
\quad\text{and}\quad
L_\alpha^{p,q}(\Omega )
 \hookrightarrow L_{\alpha_0}^{p_0,q }(\Omega )   .
\]
If $\alpha>0 > \alpha_0$, then choosing $p_1$ such that $\alpha=n/p-n/p_1$,  we have
$$
L_\alpha^{p,q}(\Omega )
 \hookrightarrow L^{p_1,q }(\Omega )
 \hookrightarrow L_{\alpha_0}^{p_0,q }(\Omega ) .
$$
This proves Part (iii).
Finally, Part (iv) is easily deduced from  Lemma \ref{interp for SL on domains} (i)  and Lemma \ref{properties of SL on Rn} (iii).
\end{proof}

\end{document}
